\chapter{From structured priorities to rational extensions}
\label{sec:p18-structured-priority-rationality}

The corrected priority constructions can be analysed directly, without
assuming the local return condition of
Definition~\ref{def:p18-rationality-return}. This chapter proves the
required return for all four choices of link principle and strict set
comparison. Shared subtrees remain in the construction.
For c-restricted arguments the return has a joint-support qualification;
a separate support lemma then proves the full extension-rationality result.

The primary sources are the corrected Modgil--Prakken general account
\cite{ModgilPrakken2013GeneralCorrected} and Caminada--Amgoud's tree
reversal lemma \cite[Lemma 6]{CaminadaAmgoud2007Evaluation}.
The latter is read in the author's 40-page file bearing the typesetting
date March 25, 2013. Its lemma retains the original assumptions.
We use that tree construction with the attack definitions already fixed
in Chapter~\ref{sec:p18-structured-argumentation}.

\section{The four direct strict priorities}

\begin{hbdefinition}{def:p18-priority-four}{Typed strict set and argument comparisons}
Fix a preorder on ordinary premises and a separate preorder on defeasible
rules. Their strict parts are denoted \(<_p\) and \(<_d\).
For either type, write \(<\) for the relevant strict part.
For nonempty finite \(\Gamma,\Delta\), define
\[
\begin{aligned}
 \Gamma\triangleleft_{\rm E}\Delta
   &\iff(\exists x\in\Gamma)(\forall y\in\Delta)\ x<y,\\
 \Gamma\triangleleft_{\rm D}\Delta
   &\iff(\forall x\in\Gamma)(\exists y\in\Delta)\ x<y.
\end{aligned}
\]
For either comparison, \(\varnothing\not\triangleleft\Delta\) for every
\(\Delta\), and \(\Gamma\triangleleft\varnothing\) for every nonempty
\(\Gamma\). The two types are never compared with one another.
These are the corrected direct strict comparisons, whose strict-order
and finite-union properties were proved in
Theorem~\ref{thm:p18-rationality-strict-sets}.

Choose \(s\in\{{\rm E},{\rm D}\}\). Last-link priority is
\[
 B\prec_{{\rm L},s}A
 \iff \operatorname{Last}(B)\triangleleft_s\operatorname{Last}(A)
 \ \lor\
 \left[
 \begin{gathered}
 \operatorname{Last}(A)=\operatorname{Last}(B)=\varnothing\\[-2pt]
 \land\ \operatorname{Prem}_p(B)\triangleleft_s\operatorname{Prem}_p(A)
 \end{gathered}\right].
\]
Weakest-link priority is defined by the following ordered cases:
\[
 B\prec_{{\rm W},s}A\iff
 \begin{cases}
 \operatorname{Prem}_p(B)\triangleleft_s\operatorname{Prem}_p(A),
       &A,B\text{ both strict},\\
 \operatorname{Def}(B)\triangleleft_s\operatorname{Def}(A),
       &A,B\text{ both firm and not both strict},\\
 \operatorname{Prem}_p(B)\triangleleft_s\operatorname{Prem}_p(A)
 \ \land\
 \operatorname{Def}(B)\triangleleft_s\operatorname{Def}(A),
       &\text{otherwise}.
 \end{cases}
\]
All comparisons used by defeat in this chapter use one of these four
strict relations.
\end{hbdefinition}

\begin{hbtheorem}{thm:p18-priority-order}{proved=annex:p18.priorityorder}{The four priorities are strict partial orders}
Each relation in Definition~\ref{def:p18-priority-four} is irreflexive
and transitive. A strict firm argument is never below another argument,
and every fallible argument is below every strict firm argument.
\end{hbtheorem}

\begin{hblemma}{thm:p18-priority-resources}{proved=annex:p18.priorityresources}{Finite resource witnesses}
Let \(\Gamma,\Delta,H\) be finite subsets of a strictly partially ordered
set.

If \(\Gamma\triangleleft_{\rm E}\Delta\) and
\(\Delta\ne\varnothing\), there is an \(m\in\Gamma\) minimal in
\(\Gamma\cup\Delta\). If \(\Delta=\varnothing\) and
\(\Gamma\ne\varnothing\), choose \(m\) minimal in \(\Gamma\).
For any finite \(\Sigma\) containing \(m\) and any
\(H\subseteq\Gamma\cup\Delta\),
\[
 H\not\triangleleft_{\rm E}\Sigma.
\]

If \(\Gamma\triangleleft_{\rm D}\Delta\) and
\(\Delta\ne\varnothing\), there is a maximal
\(d\in\Gamma\cup\Delta\) belonging to \(\Delta\).
If \(d\in H\) and \(\varnothing\ne\Sigma\subseteq\Gamma\), then
\(H\not\triangleleft_{\rm D}\Sigma\).
Also, for any nonempty finite \(\Sigma\subseteq H\),
\(H\not\triangleleft_{\rm D}\Sigma\).
\end{hblemma}

\section{Continuations and repeated indexed families}

\begin{hbdefinition}{def:p18-priority-tree-continuation}{Actual strict continuations}
For a finite indexed family \(A_1,\ldots,A_m\), a strict tree continuation
is obtained by taking copies of every indexed input tree, adding axiom
leaves if needed, and joining these trees by strict rule applications
only. Every indexed input occurs at least once. Equal trees at different
indices may have distinct occurrences. For an empty family this produces
only strict firm arguments, when such arguments exist.
This construction preserves the union of ordinary-premise, defeasible-rule,
and last-rule sets of the indexed inputs.
\end{hbdefinition}

\begin{hbtheorem}{thm:p18-priority-continuations}{proved=annex:p18.prioritycontinuations}{Reasonable priorities for actual continuations}
All four corrected priorities satisfy clauses (i)--(iv) of
Definition~\ref{def:p18-rationality-reasonable} when continuation means
strict tree continuation. Clause (iv) may quantify over arbitrary finite
indexed families, including repetitions.
For weakest link, both comparisons also satisfy those clauses for the
source's broader profile-continuation relation, including that
repetition-permitting interpretation.
\end{hbtheorem}

\begin{hbtheorem}{thm:p18-priority-profile-boundary}{proved=annex:p18.priorityprofile}{Last link distinguishes repeated defeasible applications}
For either corrected set comparison, last-link priority need not satisfy
the singleton-continuation clause (iii) when continuation means only
the component-set relation of
Definition~\ref{def:p18-rationality-reasonable}.
A theory with one axiom and two defeasible rules witnesses this failure.
Its violating continuation adds another application of an already used
defeasible rule, and is not an actual strict tree continuation.
\end{hbtheorem}

\section{Reversing strict paths with shared subtrees}

Transposition below is positionwise: replace one selected antecedent by
the contradictory of the consequent, and make the contradictory of that
antecedent the new consequent. The proof keeps the other children in their
original positions, so no separate closure under antecedent permutations
is required.

\begin{hbtheorem}{thm:p18-priority-retained-reversal}{proved=annex:p18.priorityreversal}{The assumption-retaining reversal}
Suppose a paired-literal strict-rule set is closed under transposition.
Let \(A\) be a finite strict derivation of \(c\), whose leaves include
assumptions \(a_1,\ldots,a_k\), with \(k\geq1\).
For each selected occurrence of an assumption \(a_i\), a strict derivation
of \(\neg a_i\) can be built with premises drawn from the original
assumptions together with \(\neg c\). Repeated occurrences of \(a_i\)
in other branches may remain.

More generally, if \(X\) concludes the contradictory of the conclusion
of an argument \(B\), a strict path from the root of \(B\) to a
frontier subtree \(U\in M(B)\) can be reversed to obtain a strict graft
\(W\) of \(X\) and off-path subtrees of \(B\) attacking \(B\) at \(U\).
The graft contains \(X\) as a subtree. Every defeater of \(W\) defeats
\(X\) or \(B\). Its premise, defeasible-rule, and last-rule sets are
contained in the corresponding unions for \(X,B\).
\end{hbtheorem}

\begin{hblemma}{thm:p18-priority-contraposition-graft}{proved=annex:p18.prioritycontra}{A contraposition graft from two opposing roots}
Under closure under contraposition, let \(X,B\) have contradictory
conclusions, and let \(U\in M(B)\).
There exists a strict graft \(W\) attacking \(B\) at \(U\), obtained
by substituting only whole copies of \(X\) and \(B\) into a strict
derivation. It may use either, both, or neither of those copies.
Every defeater of \(W\) defeats \(X\) or \(B\).
If it uses neither, it is strict and firm.
If c-restricted arguments are used and
\(\operatorname{Prem}(X)\cup\operatorname{Prem}(B)\) is c-consistent,
the graft is admitted.
\end{hblemma}

\begin{hbtheorem}{thm:p18-priority-return}{proved=annex:p18.priorityreturn}{Return for the corrected priorities}
Use any of the four corrected priorities. Assume either closure under
contraposition, or paired literals and closure under transposition.
Suppose \(B T A\) but not \(B D A\).
Then the attack is contradictory undermining or rebutting at a fallible
subargument \(X\) of \(A\), and \(B\prec X\).

In the unrestricted argument universe there is an argument \(W\) with
\[
 W D B,\qquad
 (\forall Z)\,[Z D W\Longrightarrow (Z D B\lor Z D A)].
 \tag{P}
\]
Thus the local return condition (R) is derived.

In the c-restricted universe the same conclusion holds whenever
\(\operatorname{Prem}(B)\cup\operatorname{Prem}(A)\) is c-consistent.
No disjointness of these premise sets or of the argument trees is
assumed. The proof permits arbitrarily many repeated occurrences of
a subargument.
\end{hbtheorem}

\section{Restricted supports and the full rationality conclusion}

\begin{hblemma}{thm:p18-priority-supported-union}{proved=annex:p18.prioritysupport}{Support consistency without assuming return}
Use a c-restricted universe with c-classical strict consequence and
axiom consistency, and any of the four corrected priorities.
If finitely many arguments are defended by a defeat-conflict-free
set \(E\), their combined premise set is c-consistent.
This result does not assume (R).
\end{hblemma}

\begin{hbtheorem}{thm:p18-priority-rationality}{proved=annex:p18.priorityrationality}{Closure and consistency under all four priorities}
Assume axiom consistency, contrary well-formedness, and either closure
under contraposition or paired-literal closure under transposition.
Use all finite arguments, or all finite arguments with c-consistent
premises; in the latter case additionally assume c-classicality.
Use any of the four corrected priority relations.

Then attack and defeat versions of admissible, complete, grounded,
preferred, and stable semantics coincide.
Every admissible extension has directly consistent conclusions.
Every complete extension contains its subarguments and is closed under
strict inference; its conclusions are indirectly consistent.
These assertions hold for infinite languages, rule sets, and argument
universes as well as finite ones.

In the unrestricted case the result follows from the proved global
return and Theorem~\ref{thm:p18-rationality-general}.
In the restricted case the proof supplies the support consistency
needed at each return step; it does not assume global (R).
\end{hbtheorem}

\begin{hbtheorem}{thm:p18-priority-restricted-boundary}{proved=annex:p18.priorityrestricted}{Why global return is too strong for restricted supports}
Even with classical strict consequence, consistent axioms,
c-classicality, and a total premise preorder, a c-restricted framework
using any of the four corrected priorities can fail global (R).
This does not contradict Theorem~\ref{thm:p18-priority-rationality}:
the failed pair in the countermodel has c-inconsistent combined support.
The proof supplies its complete rule specification and proves that
no returning defeater exists.
\end{hbtheorem}

\section{A worked priority comparison}

Take ordinary premises \(a,b\), no axioms or strict rules, and the
defeasible rules
\[
 d_1:a\Rightarrow u,\qquad
 d_2:u\Rightarrow t,\qquad
 d_3:b\Rightarrow\neg t.
\]
Let the language contain the paired literals over
\(a,b,u,t,\delta_1,\delta_2,\delta_3\), with \(n(d_i)=\delta_i\).
The five arguments are
\[
 A_0=[a],\quad B_0=[b],\quad
 I=[A_0;d_1],\quad A=[I;d_2],\quad B=[B_0;d_3].
\]
Set \(a<_p b\) and \(d_1<_d d_3<_d d_2\).

\begin{hbtheorem}{thm:p18-priority-worked}{proved=annex:p18.priorityworked}{The four comparisons and their extensions}
The only attacks are \(A T B\) and \(B T A\).
For either last-link comparison, \(B\prec A\), and the unique
complete, grounded, preferred, and stable extension is
\(\{A_0,B_0,I,A\}\).
For weakest-link elitism, \(A\prec B\), and the unique extension
in all four semantics is \(\{A_0,B_0,I,B\}\).
For weakest-link democracy neither strict comparison holds.
There are two preferred and stable extensions, those same two
four-argument sets; the complete extensions are these two sets and
\(\{A_0,B_0,I\}\), which is grounded.
Every failed attack in this example has its explicit return from the
opposing top argument.
\end{hbtheorem}

\begin{proofboundary}{def:p18-priority-next}{What the priority proof establishes}
The implication from each corrected priority construction to the
stated rationality results is proved here, including the restricted-support
case. The original Caminada--Amgoud assumption-retaining lemma is also
supplied with its complete construction.
No resource-deletion assertion or disjoint-tree restriction is needed.

The reasonable-order claim is proved for actual strict continuations
with repetitions, and also for weakest link with profile continuations.
The last-link profile claim fails by the explicit two-rule countermodel.
These distinctions do not reduce the scope of the direct return proof.
The classical/Tarskian representation theorems, preferred-subtheory
equivalence, and other
instantiation results of the general account remain separate source
presentations.
\end{proofboundary}
